Kurs:Algebraische Kurven (Osnabrück 2012)/Arbeitsblatt 24/latex
Kurs:Algebraische Kurven (Osnabrück 2012)/Arbeitsblatt 24/latex


\setcounter{section}{24}


  \zwischenueberschrift{Aufwärmaufgaben}


\inputaufgabe

{}

{


Beschreibe ein Beispiel einer glatten Kurve


\mavergleichskette

{\vergleichskette

{ C
}

{ \subseteq }{ {\mathbb A}^{2}_{K}
}

{  }{
}

{    }{
}

{   }{
}

}
{}{}{}
mit einer Parametrisierung, deren Differential an mindestens einem Punkt verschwindet.


}

{} {}


\inputaufgabe

{}

{


Betrachte das Achsenkreuz


\mavergleichskette

{\vergleichskette

{ V(xy)
}

{ \subseteq }{ {\mathbb A}^{2}_{K}
}

{  }{
}

{    }{
}

{   }{
}

}
{}{}{}
und den zum Nullpunkt gehörigen
\definitionsverweis {lokalen Ring}{}{}
$R$ mit
\definitionsverweis {maximalem Ideal}{}{}
${\mathfrak m}$. Beschreibe explizit eine
$K$-\definitionsverweis {Basis}{}{}
für die
\definitionsverweis {Restklassenringe}{}{}
$R/{\mathfrak m}^n$ und bestimme die
\definitionsverweis {Dimensionen}{}{}
davon.


}

{} {}


\inputaufgabe

{}

{


Es sei $K$ ein
\definitionsverweis {Körper}{}{}
und

\mathl{K [ \![ T]\! ]}{} der
\definitionsverweis {Potenzreihenring}{}{.}
Man gebe die inverse Potenzreihe zu

\mathl{1-T}{} an.


}

{} {}


\inputaufgabegibtloesung

{}

{


Es sei $K$ ein
\definitionsverweis {Körper}{}{,}


\mavergleichskette

{\vergleichskette

{ {\mathfrak m}
}

{ = }{ (T)
}

{ \subseteq }{ K[T]
}

{    }{
}

{   }{
}

}
{}{}{}
das zum Nullpunkt gehörige
\definitionsverweis {maximale Ideal}{}{}
mit der
\definitionsverweis {Lokalisierung}{}{}


\mavergleichskette

{\vergleichskette

{ R
}

{ = }{ K[T]_{\mathfrak m}
}

{  }{
}

{    }{
}

{   }{
}

}
{}{}{.}
Definiere einen
$K$-\definitionsverweis {Algebrahomo\-morphismus}{}{}
\maabbdisp {\varphi} { R } { K [ \![ T]\! ]
} {}
mit


\mavergleichskette

{\vergleichskette

{ \varphi(T)
}

{ = }{ T
}

{  }{
}

{    }{
}

{   }{}

}
{}{}{,}
wobei

\mathl{K [ \![ T]\! ]}{} den
\definitionsverweis {Ring der formalen Potenzreihen}{}{}
bezeichnet.


}

{} {}


\inputaufgabe

{}

{


Berechne die ersten fünf Glieder (bis einschließlich $c_4$) der eingesetzten Potenzreihe $F(G)$ im Sinne von
Definition 24.9.


}

{} {}


  \zwischenueberschrift{Aufgaben zum Abgeben}


\inputaufgabe

{}

{


Man gebe ein Beispiel für ein irreduzibles reelles Polynom


\mavergleichskette

{\vergleichskette

{ F
}

{ \in }{ \R[X,Y]
}

{  }{
}

{    }{
}

{   }{
}

}
{}{}{}
derart, dass beide partiellen Ableitungen übereinstimmen und nicht konstant sind. Zeige, dass dies über ${\mathbb C}$ nicht möglich ist.


}

{} {}


\inputaufgabe

{}

{


Betrachte die Kurve


\mavergleichskette

{\vergleichskette

{ C
}

{ = }{ V \left(  X^2-Y^2-Y^3  \right)
}

{  }{
}

{    }{
}

{   }{
}

}
{}{}{}
mit der in
Beispiel 24.3
besprochenen Parametrisierung. Bestimme die singulären Punkte der Kurve zusammen mit den Multiplizitäten und Tangenten. Berechne ebenfalls die Bildpunkte und die Tangenten für die Parameterwerte


\mavergleichskette

{\vergleichskette

{ t
}

{ = }{ -1,0,1
}

{  }{
}

{    }{
}

{   }{
}

}
{}{}{.}


}

{} {}


\inputaufgabe

{}

{


Beschreibe eine formale Potenzreihe über ${\mathbb C}$, die in keiner Umgebung des Nullpunktes konvergiert.


}

{} {}


\inputaufgabe

{}

{


Es sei $K$ ein Körper. Vergleiche die beiden Ringe
\mathkor {} {(K[X])[ \![Y]\! ]} {und} {(K [ \![ Y]\! ])[X]} {.}


}

{} {}


\inputaufgabe

{}

{


Es sei $R$ ein
\definitionsverweis {noetherscher}{}{}
kommutativer Ring. Man zeige, dass

\mathl{R [\![ T_1, \ldots ,  T_{ n }]\!]}{} noethersch ist.


}

{} {Hinweis: Lassen Sie sich vom Beweis des Hilbertschen Basissatzes inspirieren!}


<< | Kurs:Algebraische Kurven (Osnabrück 2012) | >>


PDF-Version dieses Arbeitsblattes


 Zur Vorlesung
(PDF)
